% 06-interpretacion.tex — Sección 6: interpretación. Responde M1, M2 y M3 con
% datos, y aplica las tres decisiones de Alix: R-29.1 (distinguir dónde ocurrió
% cada OOM), R-29.2 (nota al pie de la celda colapsada) y R-16.1 (la premisa del
% manual sobre las variantes de Q2, con la cobertura parcial declarada).

\section{Interpretación}
\label{sec:interpretacion}

\subsection{M1: ¿cómo se degrada cada motor al apretar los recursos?}
\label{sec:m1}

El resultado no se parte por motor, se parte \textbf{por consulta}, y es
consistente en los tres perfiles y en las tres concurrencias.

\begin{table}[htbp]
  \centering
  \caption{Quién gana cada celda, en throughput (media de 3 repeticiones). El
           factor es cuántas veces más consultas por segundo completa el ganador.}
  \label{tab:m1}
  \small
  \begin{tabular}{llrrr}
    \toprule
    Consulta & Perfil & c=1 & c=10 & c=50 \\
    \midrule
    \texttt{Q1} & A & Mongo $\times$2.2 & Mongo $\times$3.8 & PostgreSQL no completa\tnotaCensurada{} \\
    \texttt{Q1} & B & Mongo $\times$2.4 & Mongo $\times$3.3 & Mongo $\times$25 \\
    \texttt{Q1} & C & Mongo $\times$3.4 & Mongo $\times$5.4 & los dos fallan\tnotaColapsada{} \\
    \texttt{Q2} & A & PostgreSQL $\times$2.6 & PostgreSQL $\times$1.2 & Mongo $\times$1.4\tnotaBajaDensidad{} \\
    \texttt{Q2} & B & PostgreSQL $\times$6.0 & PostgreSQL $\times$2.8 & PostgreSQL $\times$2.8 \\
    \texttt{Q2} & C & PostgreSQL $\times$10.7 & PostgreSQL $\times$6.0 & PostgreSQL $\times$3.4 \\
    \bottomrule
  \end{tabular}
\end{table}

\paragraph{\texttt{Q1} (desanidar y agrupar): MongoDB gana siempre.}
Y la ventaja \textbf{crece} al apretar los recursos: de $\times$2.2 en el perfil A
a $\times$3.4 en el C con un cliente. La razón está en el plan y en el disco: para
PostgreSQL, \texttt{Q1} exige un ordenamiento que no cabe en \texttt{work\_mem} y
termina en disco \textbf{en los tres perfiles} (sección~\ref{sec:spill}), mientras
MongoDB la resuelve sin tocar disco.

\paragraph{\texttt{Q2} (ordenar y quedarse con el Top-1000): PostgreSQL gana, y la ventaja crece al apretar.}
De $\times$2.6 a $\times$10.7 con un cliente al pasar del perfil A al C. Visto al
revés: entre el perfil A y el C, MongoDB se vuelve \textbf{13.7 veces} más lento
en \texttt{Q2} y solo 1.5 veces más lento en \texttt{Q1}. La sensibilidad a la
contención depende más de la consulta que del motor.

\paragraph{La única excepción, y por qué no se usa.}
La celda \texttt{Q2} · perfil~A · \mbox{c=50} es la única de \texttt{Q2} que gana
MongoDB, y es también \textbf{una de las dos celdas con dispersión por encima del
15~\%} (sección~\ref{sec:estabilidad}): su medición de PostgreSQL cayó un 34~\%
entre la primera y la tercera repetición mientras el procesador bajaba de
frecuencia. Esa inversión no se usa para sostener nada.

\subsection{M2: ¿hubo volcado a disco y OOM, y dónde?}
\label{sec:spill}

\paragraph{Volcado: sí, masivo, y solo en PostgreSQL con \texttt{Q1}.}
\cifraFilasConSpill{} de las \cifraFilasCsv{} filas tienen
\texttt{disk\_spill = true}, y \textbf{todas son de PostgreSQL con \texttt{Q1}}:
8 de sus 9 celdas (la novena no volcó porque el motor murió antes). En total,
\textbf{\cifraTempTotalTB{}~TB} de archivos temporales escritos a lo largo de la
matriz. Con un solo cliente en el perfil A son unos 32~GB por ventana de 300~s,
es decir \textbf{unos 0.8~GB por consulta}.

\textbf{MongoDB no volcó en ninguna celda}, ni siquiera con
\texttt{allowDiskUse: true} (anexo~\ref{sec:variantes}).

\paragraph{OOM: cinco filas, y hay que distinguir dónde ocurrió cada uno.}
Son \cifraFilasOom{} filas con \texttt{status = oom}, todas de PostgreSQL con
\texttt{Q1} y \mbox{c=50} en los perfiles B y C. Presentarlas como «cinco OOM» sin
más sería impreciso:

\begin{itemize}
  \item \textbf{\cifraOomEnVentana{} ocurrieron dentro de la ventana de medición}:
        el contador \texttt{oom\_kill} del cgroup se movió mientras se medía.
  \item \textbf{\cifraOomEnCalentamiento{} ocurrieron durante el calentamiento}: el
        contador ya estaba en su valor final al abrirse la ventana. El motor llegó
        a la medición \emph{ya} sin memoria.
\end{itemize}

La distinción importa porque cambia lo que la fila mide. Las tres de la segunda
clase no describen «un motor que se queda sin memoria atendiendo la carga medida»,
sino «un motor que ya se había quedado sin memoria antes de empezar a medir». El
detalle de cada una está en \texttt{results/logs/}.

\begin{table}[htbp]
  \centering
  \caption{Las \cifraFilasOom{} repeticiones con \texttt{status = oom}. La columna
           «$\Delta$ \texttt{oom\_kill}» es cuántos procesos mató el kernel
           \emph{durante la ventana}: un cero con \texttt{OOMKilled = sí}
           significa que el OOM ocurrió antes, en el calentamiento (R-29.1).}
  \label{tab:oom}
  \small
  \begin{tabular}{lllrrrlrl}
    \toprule
    Motor & Perfil & Consulta & c & Rep. & $\Delta$ \texttt{oom\_kill} & Cuándo &
    Completadas & \texttt{OOMKilled} \\
    \midrule
    % generado por scripts/build_report_data.py — no editar a mano
PostgreSQL & B & Q1 & 50 & 1 & 0 & en el calentamiento & 4 & sí \\
PostgreSQL & B & Q1 & 50 & 2 & 0 & en el calentamiento & 4 & sí \\
PostgreSQL & C & Q1 & 50 & 1 & 0 & en el calentamiento & 0 & sí \\
PostgreSQL & C & Q1 & 50 & 2 & 1 & en la ventana & 0 & sí \\
PostgreSQL & C & Q1 & 50 & 3 & 6 & en la ventana & 0 & sí \\
   % el archivo generado cierra con \bottomrule (ver report/README.md)
  \end{tabular}
\end{table}

\paragraph{La peor celda, y cómo no leerla.}
En \textbf{PostgreSQL · perfil C · \texttt{Q1} · \mbox{c=50}} no se completó
\emph{ninguna} consulta en las tres repeticiones: el throughput es 0.000. La fila
registra \texttt{cpu\_throttled\_seconds = 0.035} y
\texttt{max\_rss\_mb = 8.2}\tnotaColapsada{}, y esos dos números \textbf{no
significan que PostgreSQL atendiera 50 clientes con 8~MiB de memoria}: significan
que durante la ventana el motor estaba colapsado y no trabajaba. Sus
\num{15000} errores registrados son \textbf{conexiones rechazadas}, no consultas
ejecutadas. Por eso esa celda lleva nota al pie en todas las tablas de este
reporte.

\paragraph{Respuesta a M2.} El volcado a disco es real, sistemático y enorme, pero
está concentrado en una sola combinación: PostgreSQL con \texttt{Q1}. El OOM
aparece solo donde se juntan la consulta más pesada, la concurrencia más alta y
los perfiles con menos memoria. MongoDB no volcó ni fue matado por el kernel en
ninguna celda.

\subsection{La premisa del manual sobre las dos variantes de \texttt{Q2}}
\label{sec:variantes}

El manual pide dos variantes de estrés de memoria para \texttt{Q2}
(\texttt{work\_mem} de 4~MB frente a 64~MB en PostgreSQL;
\texttt{allowDiskUse} falso frente a verdadero en MongoDB) \textbf{dando por
supuesto que se diferenciarían}. Se midieron las dos en el perfil con menos
memoria y \textbf{no se diferencian}: ni en throughput más allá de la dispersión
propia de cada celda, ni en volcado, que no aparece en ninguna de las dos.

La explicación está en el plan (sección~\ref{sec:planes}): el
\texttt{LIMIT 1000} convierte el ordenamiento en un \emph{top-K} de
\textbf{273~kB}, que caben de sobra en los 4~MB de la variante restrictiva. El
parámetro que el manual varía nunca llega a ser el límite. El detalle de la
comparación está en el anexo~\ref{sec:variantes-anexo}, con la cobertura parcial
declarada.

\subsection{M3: huella de CPU y memoria por cada 1000 documentos}
\label{sec:memoria}

El CSV no tiene un denominador, así que se toma de los \textbf{planes de
ejecución}, como corresponde: los documentos que cada consulta examina de verdad.
Y ahí aparece una asimetría que hay que declarar antes de dividir:

\begin{center}
  \small
  \begin{tabular}{llr}
    \toprule
    Motor & Consulta & Documentos examinados por consulta \\
    \midrule
    MongoDB    & \texttt{Q1} & \num{500000} (recorrido completo) \\
    MongoDB    & \texttt{Q2} & \num{250207} (usa su índice de \texttt{timestamp}) \\
    PostgreSQL & \texttt{Q1} & \num{500000} \\
    PostgreSQL & \texttt{Q2} & \num{500000} (\texttt{Seq Scan}: no usa el índice) \\
    \bottomrule
  \end{tabular}
\end{center}

\paragraph{CPU por 1000 documentos (perfil A, un cliente).}
MongoDB gasta 8.6~ms en \texttt{Q1} y PostgreSQL 25.6~ms: \textbf{MongoDB es unas
tres veces más eficiente en CPU} para esa consulta. En \texttt{Q2} se invierte:
PostgreSQL gasta 3.5~ms y MongoDB 6.6~ms. Y al apretar los recursos la brecha de
\texttt{Q2} se abre hasta ser de \textbf{74.7~ms frente a 3.6~ms} en el perfil~C,
unas 21 veces.

\paragraph{La memoria depende de cuál de las dos memorias se mire, y esto es importante.}
\texttt{max\_rss\_mb} mide la memoria \emph{anónima} del contenedor. Los
\texttt{shared\_buffers} de PostgreSQL son memoria \textbf{compartida}, que el
kernel no cuenta como anónima; el caché de WiredTiger de MongoDB sí lo es. El
efecto es grande:

\begin{center}
  \small
  \begin{tabular}{lrrrr}
    \toprule
    & \multicolumn{2}{c}{\texttt{max\_rss\_mb} (anónima)} & \multicolumn{2}{c}{\texttt{working\_set}} \\
    \cmidrule(lr){2-3}\cmidrule(lr){4-5}
    Perfil A, c=1 & \texttt{Q1} & \texttt{Q2} & \texttt{Q1} & \texttt{Q2} \\
    \midrule
    MongoDB    & 824 & 832 & 1042 & 1058 \\
    PostgreSQL & 104 & 16  & 1630 & 891 \\
    \bottomrule
  \end{tabular}
\end{center}

Leyendo solo \texttt{max\_rss\_mb}, PostgreSQL parecería de 8 a 50 veces más
eficiente en memoria. Leyendo \texttt{working\_set} —que sí incluye la memoria
compartida y es la que cuenta contra el límite del contenedor— \textbf{la relación
se invierte en \texttt{Q1}}: 1630~MiB de PostgreSQL frente a 1042 de MongoDB. Por
eso M3 se responde con \texttt{working\_set}, y la figura~\ref{fig:curva-rss}
muestra las dos curvas juntas.

\paragraph{Un matiz honesto sobre \texttt{working\_set}.}
Incluye también la caché de páginas de los archivos de datos que se le atribuye al
contenedor. PostgreSQL depende de ella por diseño —lee a través de la caché del
sistema operativo— y cuenta contra su \texttt{memory.max}, así que es memoria que
el motor necesita bajo el límite; pero no es «memoria del proceso». Las dos
lecturas están en los datos (\texttt{report/datos/m3\_huella.dat}) para que
cualquiera pueda rehacer la cuenta con el criterio que prefiera.

\subsection{Y la conclusión que el throttling deja clara}

El throttling confirma que los perfiles del manual \textbf{aprietan de verdad}:
hasta \num{2086}~s acumulados en una ventana de 300~s. Pero no es lo que
diferencia a los motores. Lo que los diferencia es si la consulta cabe en memoria:
donde cabe (\texttt{Q2}, por el \emph{top-K}), gana el motor con el plan más
eficiente; donde no cabe (\texttt{Q1} en PostgreSQL), el disco domina y la CPU
pasa a segundo plano.
